\paragrafos{Ο κύκλος και τα στοιχεία του}{Β.1.11, «Κύκλος και στοιχεία του κύκλου»}

\begin{anagke}
Το πιο όμορφο σημείο του κήπου θα είναι ένα στρογγυλό παρτέρι με ένα συντριβάνι στο κέντρο. Ο παππούς Γιάννης δείχνει πώς θα το χαράξουν στο χώμα: καρφώνει ένα παλούκι στο κέντρο, δένει ένα σχοινί και, κρατώντας το τεντωμένο, γυρίζει και χαράσσει τη γραμμή. «Όλα τα σημεία της γραμμής έχουν την ίδια απόσταση από το παλούκι, όσο είναι το σχοινί», λέει.

Η Αγγελική το αναγνωρίζει: «Αυτό είναι ο κύκλος! Και το σχοινί είναι η ακτίνα!». Ο Νίκος ρωτάει πώς λέγεται το πιο μακρύ τμήμα που μπορούν να τεντώσουν από ένα άκρο του κύκλου ως το άλλο: «Η διάμετρος», απαντά η Πολυξένη, «περνάει από το κέντρο και είναι διπλάσια από την ακτίνα». Ο Βασίλης στέκεται κάπου και ρωτάει αν είναι «μέσα» ή «έξω» από τον κύκλο: «Αν η απόστασή σου από το παλούκι είναι μικρότερη από το σχοινί, είσαι μέσα».

Ο κύκλος είναι το πιο συμμετρικό σχήμα και όλα τα στοιχεία του προκύπτουν από τη μέτρηση αποστάσεων από το κέντρο.
\end{anagke}

\begin{lexi}{κύκλος, ακτίνα, διάμετρος, χορδή}{κύκλος, ακτίνα, διάμετρος, χορδή}
Ο \lexiword{κύκλος} είναι ό,τι «κυλάει», γυρίζει. Η \lexiword{ακτίνα} είναι η «αχτίδα», η ακτίνα του φωτός που ξεκινά από το κέντρο. Η \lexiword{διάμετρος} είναι το «δια-μέτρο», αυτό που μετράει «δια μέσου» του κέντρου. Η \lexiword{χορδή} είναι η «χορδή» του μουσικού οργάνου: ένα τεντωμένο νήμα ανάμεσα σε δύο σημεία. Το \lexiword{ημικύκλιο} είναι το μισό του κύκλου.
\end{lexi}

\begin{metafrasi}
Ο κύκλος με κέντρο $\mathrm{Ο}$ και ακτίνα $\rho$. & $(\mathrm{Ο},\rho)$\\
Τα σημεία πάνω στον κύκλο. & $\mathrm{ΟΜ}=\rho$\\
Τα σημεία μέσα στον κύκλο. & $\mathrm{ΟΜ}<\rho$\\
Τα σημεία έξω από τον κύκλο. & $\mathrm{ΟΜ}>\rho$\\
Η διάμετρος. & $\delta=2\rho$\\
\end{metafrasi}

\tmima{Η θεωρία}

\begin{orismos}{κύκλος}
\textbf{Κύκλος} με κέντρο $\mathrm{Ο}$ και ακτίνα $\rho$ είναι το σύνολο όλων των σημείων του επιπέδου που απέχουν από το $\mathrm{Ο}$ απόσταση ίση με $\rho$. Συμβολίζεται $(\mathrm{Ο},\rho)$. Η απόσταση $\rho$ λέγεται \textbf{ακτίνα} και το $\mathrm{Ο}$ \textbf{κέντρο}.
\end{orismos}

Για να σχεδιάσουμε κύκλο χρησιμοποιούμε \textbf{διαβήτη}. Δύο κύκλοι με ίσες ακτίνες είναι ίσοι.

\begin{center}
\begin{tikzpicture}[scale=1.15]
\draw[thick] (0,0) circle (1.6);
\fill (0,0) circle (2pt) node[below left] {$\mathrm{Ο}$};
\draw[gblue,thick] (0,0)--(1.13,1.13) node[midway,above left,gblue] {$\rho$}; \fill (1.13,1.13) circle (2pt) node[above right] {$\mathrm{Α}$};
\draw (-1.6,0)--(1.6,0); \fill (-1.6,0) circle (2pt) node[left] {$\mathrm{Β}$}; \fill (1.6,0) circle (2pt) node[right] {$\mathrm{Γ}$};
\node[below] at (0,-0.1) {\small διάμετρος $\mathrm{ΒΓ}$};
\draw[dashed] (-1.2,1.06)--(0.8,-1.38); \fill (-1.2,1.06) circle (2pt); \fill (0.8,-1.38) circle (2pt);
\node at (-1.2,-0.9) {\small χορδή};
\end{tikzpicture}
\end{center}

\begin{itemize}[leftmargin=1.2em,itemsep=0.3ex,topsep=2pt]
\item \textbf{Χορδή} είναι το ευθύγραμμο τμήμα που ενώνει δύο σημεία του κύκλου.
\item \textbf{Διάμετρος} είναι η χορδή που περνά από το κέντρο. Είναι η μεγαλύτερη χορδή, ίση με $\delta=2\rho$, και χωρίζει τον κύκλο σε δύο ίσα μέρη, τα \textbf{ημικύκλια}.
\item Ένα σημείο $\mathrm{Μ}$ είναι \textbf{πάνω} στον κύκλο αν $\mathrm{ΟΜ}=\rho$, \textbf{μέσα} σε αυτόν αν $\mathrm{ΟΜ}<\rho$ και \textbf{έξω} αν $\mathrm{ΟΜ}>\rho$.
\item Ο κύκλος μαζί με το εσωτερικό του λέγεται \textbf{κυκλικός δίσκος}.
\end{itemize}

\begin{prosoxi}
\begin{itemize}[leftmargin=1.2em,itemsep=0.3ex,topsep=2pt]
\item Ο κύκλος είναι μόνο η γραμμή, όχι το εσωτερικό: το εσωτερικό μαζί με τη γραμμή είναι ο κυκλικός δίσκος.
\item Η ακτίνα είναι τμήμα (ή μήκος), όχι ολόκληρη ευθεία.
\item Η διάμετρος είναι πάντα $2\rho$, και κάθε διάμετρος είναι χορδή, αλλά όχι κάθε χορδή είναι διάμετρος.
\item Σε έναν κύκλο όλες οι ακτίνες είναι ίσες, όλες οι διάμετροι είναι ίσες.
\end{itemize}
\end{prosoxi}

\tmima{Λυμένα παραδείγματα}

\paradeigma{Το κυκλικό παρτέρι}
\tixero Το σχοινί του παππού είναι $2{,}5$ μ.\par
\tithelo Η ακτίνα και η διάμετρος του παρτεριού.\par
\possindeo Το σχοινί είναι η ακτίνα: $\rho=2{,}5$ μ. Η διάμετρος $\delta=2\cdot2{,}5=5$ μ.\par
\apantisi $\rho=2{,}5$ μ. και $\delta=5$ μ.\par
\elegxos Το παρτέρι έχει πλάτος $5$ μ.

\paradeigma{Μέσα, πάνω ή έξω;}
\tixero Κύκλος $(\mathrm{Ο},4)$ και σημεία με $\mathrm{ΟΑ}=3$, $\mathrm{ΟΒ}=4$, $\mathrm{ΟΓ}=5$.\par
\tithelo Η θέση κάθε σημείου.\par
\possindeo Συγκρίνω την απόσταση με την ακτίνα $4$. $\mathrm{ΟΑ}<4$: μέσα. $\mathrm{ΟΒ}=4$: πάνω. $\mathrm{ΟΓ}>4$: έξω.\par
\apantisi Το $\mathrm{Α}$ μέσα, το $\mathrm{Β}$ πάνω, το $\mathrm{Γ}$ έξω.\par
\elegxos Η απόσταση από το κέντρο καθορίζει τη θέση.

\paradeigma{Μία χορδή}
\tixero Ο κύκλος $(\mathrm{Ο},5)$ και μια χορδή $\mathrm{ΑΒ}$.\par
\tithelo Πόσο μεγάλη μπορεί να είναι η χορδή.\par
\possindeo Η μεγαλύτερη χορδή είναι η διάμετρος, $2\cdot5=10$.\par
\apantisi $\mathrm{ΑΒ}\le10$, ίσο με $10$ μόνο όταν περνά από το κέντρο.\par
\elegxos Αν $\mathrm{ΑΒ}=10$ τότε το κέντρο είναι το μέσο του $\mathrm{ΑΒ}$.

\paradeigma{Ίσοι κύκλοι}
\tixero Δύο κύκλοι: $(\mathrm{Ο},3)$ και $(\mathrm{Κ},3)$.\par
\tithelo Είναι ίσοι;\par
\possindeo Έχουν ίσες ακτίνες, άρα μπορούν να εφαρμοστούν ο ένας πάνω στον άλλο μετακινώντας το κέντρο.\par
\apantisi Ναι.\par
\elegxos Ίσες ακτίνες δίνουν ίσους κύκλους.

\tmima{Ασκήσεις}

\zoni{Το κατάλαβα}{ακτίνα, διάμετρος, χορδή}

\ex Σχεδίασε με διαβήτη κύκλο $(\mathrm{Ο},3\ \mathrm{cm})$. Φέρε μια ακτίνα, μια διάμετρο και μια χορδή που δεν περνά από το κέντρο.

\ex Συμπλήρωσε: α) $\rho=4$, $\delta=\ldots$ \quad β) $\delta=9$, $\rho=\ldots$ \quad γ) $\rho=2{,}5$, $\delta=\ldots$ \quad δ) $\delta=7{,}2$, $\rho=\ldots$.

\ex Κύκλος $(\mathrm{Ο},4)$. Πες αν είναι μέσα, πάνω ή έξω τα σημεία με $\mathrm{ΟΑ}=3$, $\mathrm{ΟΒ}=4$, $\mathrm{ΟΓ}=5$, $\mathrm{ΟΔ}=0$.

\ex Στον κύκλο $(\mathrm{Ο},5)$ πόσο μεγάλη είναι η μεγαλύτερη χορδή; Πόσες διάμετροι υπάρχουν;

\ex Σχεδίασε δύο κύκλους με την ίδια ακτίνα και δείξε ότι είναι ίσοι με διαφανές χαρτί.

\ex Σωστό ή λάθος; α) Η διάμετρος είναι χορδή. \quad β) Η μεγαλύτερη χορδή είναι η ακτίνα. \quad γ) $\delta=2\rho$. \quad δ) Ένα σημείο με $\mathrm{ΟΜ}<\rho$ είναι πάνω στον κύκλο. \quad ε) Όλες οι ακτίνες ενός κύκλου είναι ίσες.

\zoni{Το χρησιμοποιώ}{παρτέρια, ρολόγια, ρόδες}

\ex \textbf{Το παρτέρι.} Το κυκλικό παρτέρι χαράσσεται με σχοινί $2{,}5$ μ. Πόσο είναι το πλάτος του;

\ex \textbf{Το ποδήλατο.} Ο τροχός ενός ποδηλάτου έχει διάμετρο $66$ cm. Πόση είναι η ακτίνα του; Σε ύψος πόσο βρίσκεται ο άξονάς του από το έδαφος;

\ex \textbf{Το ρολόι.} Ο λεπτοδείκτης έχει μήκος $8$ cm. Σε ποιο κύκλο κινείται η άκρη του και πόση διάμετρο έχει;

\ex \textbf{Η σέντρα.} Ο κύκλος της σέντρας σε ένα γήπεδο ποδοσφαίρου έχει ακτίνα $9{,}15$ μ. Πόση είναι η διάμετρός του;

\ex \textbf{Η στρογγυλή πίτα.} Μια πίτα έχει διάμετρο $26$ cm. Θα χωρέσει σε ταψί διαστάσεων $25\times25$ cm; Εξήγησε.

\zoni{Το σκέφτομαι}{όλα ίσα από το κέντρο}

\ex Γιατί η διάμετρος είναι η μεγαλύτερη χορδή; (Σκέψου το τμήμα ανάμεσα σε δύο σημεία και την απόστασή τους από το κέντρο.)

\ex Αν σε έναν κύκλο οι ακτίνες είναι $5$ cm, πόσα σημεία του κύκλου απέχουν $5$ cm από ένα δοσμένο σημείο του κύκλου; Πρόβλεψε και σχεδίασε.

\ex Πόσοι κύκλοι μπορούν να περάσουν από δύο δοσμένα σημεία $\mathrm{Α}$ και $\mathrm{Β}$; Πού βρίσκονται τα κέντρα τους;

\ex Μπορούν δύο διαφορετικοί κύκλοι με το ίδιο κέντρο να έχουν κοινό σημείο; Εξήγησε.

\zoni{Για το σπίτι}{στρογγυλά αντικείμενα}

\ex Βρες πέντε κυκλικά αντικείμενα στο σπίτι και μέτρησε διάμετρο και ακτίνα.

\ex Σχεδίασε έναν κύκλο με νήμα και καρφίτσα στο χαρτόνι.
